\section{Bölüm sonu çözümlü problemler}

\begin{enumerate}
  \item \textbf{Ortalama için güven aralığı.} $n=25$, $\bar x=60$, $s=10$ ve $t^*=2{,}064$ olduğunda yüzde 95 güven aralığını bulun.

  \textbf{Çözüm:} $SE=10/\sqrt{25}=2$ ve hata payı $2{,}064(2)=4{,}128$'dir. Aralık $[55{,}87;64{,}13]$ olur. Bu, bireysel puanların yüzde 95'inin bulunduğu aralık değildir.

  \item \textbf{İki ortalamanın farkı.} Tahmini fark 5, standart hata 2 ve yaklaşık kritik değer 2,01 olsun. Yüzde 95 güven aralığını hesaplayıp yorumlayın.

  \textbf{Çözüm:} $5\pm2{,}01(2)=5\pm4{,}02$; aralık $[0{,}98;9{,}02]$'dir. Veri, birinci grubun evren ortalamasının yaklaşık 1 ile 9 puan daha yüksek olduğu değerlerle uyumludur; sıfır aralıkta değildir.

  \item \textbf{Hassasiyet planlama.} Aynı değişkenlik ve güven düzeyinde hata payını 6'dan 3'e düşürmek için örneklem hacmi yaklaşık kaç katına çıkarılmalıdır?

  \textbf{Çözüm:} Hata payı yaklaşık $1/\sqrt n$ ile değişir. Hata payını yarıya indirmek için $\sqrt n$ iki kat, dolayısıyla $n$ dört kat olmalıdır. Bu ilişki yanlılık ve tasarım sorunlarını çözmez.
\end{enumerate}